\documentclass[11pt]{article}
\usepackage[T1]{fontenc}
\usepackage{textcomp}
\usepackage[margin=0.75in]{geometry}
\usepackage{amsmath,amssymb,amsthm}
\usepackage{array,booktabs}
\usepackage{hyperref}
\setlength{\parindent}{0pt}
\newtheorem{theorem}{Theorem}
\theoremstyle{definition}
\newtheorem{definition}{Definition}
\title{Flow Proof Artifact: inscribed-angle-half-central}
\author{Generated by \texttt{flow doc proof}}
\date{Flow Proof Book}
\begin{document}
\maketitle
An inscribed angle equals half the central angle subtending the same arc.\\[0.5em]
\textbf{Source.} euclid — Elements, Book III, Proposition 20\\[0.5em]
\bigskip
\noindent{\large \textbf{Derived fact 1.} \textit{An inscribed angle equals half the central angle subtending the same arc} \hfill the Euclidean plane \cdot circle \cdot inscribed angle is half the central angle}\par
\smallskip\noindent\textit{Coordinate: the Euclidean plane \cdot circle \cdot inscribed angle is half the central angle}\par
\smallskip\noindent\textit{Source: euclid — Elements, Book III, Proposition 20}\par
\smallskip\noindent\textit{Built on: all radii from the centre to the circumference are equal}\par
\medskip
\noindent\textbf{Goal.} An inscribed angle equals half the central angle subtending the same arc\par
\noindent$\displaystyle \angle APB = \tfrac{1}{2}\angle AOB$\par
\medskip
\noindent
\begin{tabular}{@{} >{\raggedright\arraybackslash}p{0.06\textwidth} >{\raggedright\arraybackslash}p{0.47\textwidth} >{\raggedleft\arraybackslash}p{0.41\textwidth} @{}}
\textbf{\#} & \textbf{Proof} & \textbf{Mathematics} \\
\midrule
\textbf{1.} & We prove that inscribed angle is half the central angle for circles in the Euclidean plane. & \\[0.45em]
\textbf{2.} & We invoke the definitional clause governing circles in the Euclidean plane: all radii from the centre to the circumference are equal. & \\[0.45em]
\textbf{3.} & From step 2, the two radii OA and OB form an isosceles triangle, so the inscribed angle at P is half the central angle at O, so angle APB equals half of angle AOB. Hence proven. & $\displaystyle \angle APB = \tfrac{1}{2}\angle AOB$ \\[0.45em]
\bottomrule
\end{tabular}
\label{thm:Geometry-circle-inscribed_angle_is_half_the_central_angle}
\par\medskip\hrule
\end{document}
